% supplementary_material_strings.tex
%
% Build with:  make strings-paper      (from the repository root)
% or:          python3 paper/make_strings_facts.py && pdflatex supplementary_material_strings.tex (x2)
%
% Every number quoted in this document is recomputed from the pyCICY source by
% paper/make_strings_facts.py and written to figures/strings_facts.tex, which
% is \input{} below. The prose and the code therefore cannot drift apart
% silently. There are no figures; the content is arithmetic.

\documentclass[11pt,a4paper]{article}

\usepackage[T1]{fontenc}
\usepackage[utf8]{inputenc}
\usepackage[margin=1in]{geometry}
\usepackage[protrusion=true,expansion=false]{microtype}
\usepackage{amsmath,amssymb,amsthm}
\usepackage{graphicx}
\usepackage{booktabs}
\usepackage{enumitem}
\usepackage{xcolor}
\usepackage[hidelinks]{hyperref}
\usepackage[nameinlink,capitalize]{cleveref}
\usepackage{framed}
\usepackage{orcidlink}

\graphicspath{{figures/}}

% Numbers computed by paper/make_strings_facts.py. Defined here as fallbacks so
% the document still typesets before the script has been run; the \input below
% overrides every one of them with the computed value.
\newcommand{\SFMatterFreeCount}{?}\newcommand{\SFMatterFreeList}{?}
\newcommand{\SFSUFiveMinusOneFund}{?}\newcommand{\SFSUFiveMinusOneAnti}{?}
\newcommand{\SFSUFiveMinusTwoFund}{?}\newcommand{\SFSUFiveZeroFund}{?}
\newcommand{\SFSUFiveZeroAnti}{?}\newcommand{\SFESevenHalf}{?}
\newcommand{\SFESevenCharged}{?}\newcommand{\SFSOTenSpinorDim}{?}
\newcommand{\SFSOTenSpinorA}{?}\newcommand{\SFSOTenMinusThreeVec}{?}
\newcommand{\SFSOTenMinusThreeSpin}{?}\newcommand{\SFSOTwelveHalfSpinor}{?}
\newcommand{\SFSOSevenIndexOne}{?}\newcommand{\SFClusterCount}{?}
\newcommand{\SFGTwoSevens}{?}\newcommand{\SFGTwoSUTwoDoublets}{?}
\newcommand{\SFGTwoSharedMult}{?}\newcommand{\SFSOSevenSpinors}{?}
\newcommand{\SFClusterThreeCharged}{?}\newcommand{\SFPTwoHfour}{?}
\newcommand{\SFPTwoHsix}{?}\newcommand{\SFPTwoChiT}{?}
\newcommand{\SFPTwoModuli}{?}\newcommand{\SFFTwoChiT}{?}
\newcommand{\SFFTwoModuli}{?}\newcommand{\SFModuliAgree}{?}
\newcommand{\SFPTwoHone}{?}\newcommand{\SFPTwoHtwo}{?}
\newcommand{\SFPTwoChi}{?}\newcommand{\SFPTwoAlg}{?}
\newcommand{\SFFZeroHone}{?}\newcommand{\SFFZeroHtwo}{?}
\newcommand{\SFFZeroChi}{?}\newcommand{\SFFZeroAlg}{?}
\newcommand{\SFFTwelveHone}{?}\newcommand{\SFFTwelveHtwo}{?}
\newcommand{\SFFTwelveChi}{?}\newcommand{\SFFTwelveAlg}{?}
\newcommand{\SFFSevenHone}{?}\newcommand{\SFFSevenHtwo}{?}
\newcommand{\SFFSevenChi}{?}\newcommand{\SFFSevenAlg}{?}
\newcommand{\SFFSevenCharged}{?}\newcommand{\SFPTwoChern}{?}
\newcommand{\SFFTwelveChern}{?}\newcommand{\SFFTwelveGap}{?}
\newcommand{\SFMWCharge}{?}\newcommand{\SFMWHone}{?}
\newcommand{\SFMWHtwo}{?}\newcommand{\SFMWGood}{?}
\newcommand{\SFPThreeCoCt}{?}\newcommand{\SFPThreeCoCube}{?}
\newcommand{\SFPThreeHone}{?}\newcommand{\SFPThreeHthree}{?}
\newcommand{\SFPThreeChi}{?}\newcommand{\SFPThreeChiChern}{?}
\newcommand{\SFPThreeTadpole}{?}\newcommand{\SFPThreeAgree}{?}
\newcommand{\SFPOneTwoCoCt}{?}\newcommand{\SFPOneTwoCoCube}{?}
\newcommand{\SFPOneTwoHone}{?}\newcommand{\SFPOneTwoHthree}{?}
\newcommand{\SFPOneTwoChi}{?}\newcommand{\SFPOneTwoChiChern}{?}
\newcommand{\SFPOneTwoTadpole}{?}\newcommand{\SFPOneTwoAgree}{?}
\newcommand{\SFPThreeFluxBound}{?}\newcommand{\SFQuinticFixedChi}{?}
\newcommand{\SFQuinticHtwoPlus}{?}\newcommand{\SFQuinticHtwoMinus}{?}
\newcommand{\SFQuinticDefPlus}{?}\newcommand{\SFQuinticDefMinus}{?}
\newcommand{\SFQuinticRoutesAgree}{?}\newcommand{\SFQuinticChiral}{?}
\newcommand{\SFQuinticVectors}{?}\newcommand{\SFQuinticFiveChiral}{?}
\newcommand{\SFQuinticFiveVectors}{?}\newcommand{\SFQuinticFiveHtwoPlus}{?}
\newcommand{\SFQuinticFiveHtwoMinus}{?}\newcommand{\SFSenCoverChi}{?}
\newcommand{\SFSenBases}{?}\newcommand{\SFSenGenusPTwo}{?}
\newcommand{\SFSenGenusDPEight}{?}\newcommand{\SFSenTadpoleAll}{?}
\newcommand{\SFBraneKodairaAgree}{?}\newcommand{\SFCicyTotal}{?}
\newcommand{\SFCicyElliptic}{?}\newcommand{\SFCicyKThree}{?}
\newcommand{\SFBicubicHone}{?}\newcommand{\SFBicubicHtwo}{?}
\newcommand{\SFBicubicFibrations}{?}\newcommand{\SFTetraFibrations}{?}
\newcommand{\SFMFsuThree}{?}\newcommand{\SFMFsoEight}{?}
\newcommand{\SFMFfFour}{?}\newcommand{\SFMFeSix}{?}
\newcommand{\SFMFeSeven}{?}\newcommand{\SFMFeEight}{?}
\IfFileExists{figures/strings_facts.tex}{\input{figures/strings_facts.tex}}{}

\newcommand{\PP}{\mathbb{P}}
\newcommand{\cO}{\mathcal{O}}
\newcommand{\ch}{\operatorname{ch}}
\newcommand{\g}[1]{\mathfrak{#1}}

\theoremstyle{definition}
\newtheorem{remark}{Remark}

\title{\vspace{-1cm}\textbf{What Anomaly Cancellation Determines,\\
and What Does Not Exist:\\[2pt]
\large Exact Methods for F-theory and Type IIB Orientifolds}}
\author{B.S. Hartshorn \orcidlink{0009-0004-2853-655X} \small \url{https://github.com/brentharts/CICY}}

\begin{document}
\maketitle

\begin{abstract}
A companion paper~\cite{SuppMath} asked which quantities in a heterotic
compactification are determined by the configuration matrix alone, and drew a
sharp line between what is exact and what is not. This paper asks the same
question of three constructions where the geometry is usually not a complete
intersection at all --- F-theory in six and four dimensions, and Type IIB
orientifolds --- and finds that the determining structure is different. It is
not the degrees of the defining polynomials. It is anomaly cancellation, and
in six dimensions it does not merely constrain the spectrum but produces it:
the matter content of a gauge algebra on a divisor is the unique solution of
three linear conditions in the intersection form of a surface, with no
cohomology and no metric anywhere. We derive rather than tabulate the six
matter-free non-Higgsable algebras, the matter of the multi-curve clusters,
the trace coefficients of the $\g{so}(N)$ spinors, and the
half-hypermultiplet of $\g{e}_7$ on a $-7$ curve. Several quantities are
computed twice by unrelated routes and required to agree --- the complex
structure moduli of a Weierstrass model by Riemann--Roch against the
gravitational anomaly, the Euler characteristic of an elliptic fourfold by a
moduli count against a Chern number, the equivariant Hodge numbers of an
orientifold involution by the Lefschetz fixed point theorem against a count
of monomials --- and each agreement is a check rather than a construction.
The boundary is again sharp, and it acquires a category the heterotic paper
had no occasion to name. Beside quantities that are exact, quantities awaiting
an implementation, and quantities obstructed by the Ricci-flat metric, there
are quantities that \emph{do not exist}: six-dimensional $(1,0)$ supersymmetry
forbids a superpotential, so an F-theory compactification to six dimensions
has no Yukawa couplings at all, and returning zero for one would be
numerically right and conceptually wrong.
\end{abstract}

\tableofcontents

% =====================================================================

\section{Introduction}
\label{sec:intro}

The companion paper~\cite{SuppMath} is organised around a single question:
given a configuration matrix, which physical quantities follow from the
degrees alone? The answer there is that a surprising number do, that they are
exact, and that the ones which do not fail for three distinguishable reasons.

That question does not transplant directly. The generic F-theory background is
an elliptic fibration whose fibre is a hypersurface in the weighted projective
space $\PP^{2,3,1}$, which is not a product of projective spaces, so it has no
configuration matrix; and an orientifold is not a manifold at all but a
manifold with an involution. What survives the move is the shape of the
question --- what is determined, by what, and exactly --- and the answer turns
out to be sharper rather than weaker.

In six dimensions the determining structure is anomaly cancellation. A chiral
theory with $\mathcal{N}=(1,0)$ supersymmetry is inconsistent unless a set of
polynomial identities in divisor classes holds~\cite{GSW1985,Sagnotti1992,
Erler1994,KumarMorrisonTaylor2010}, and those identities are strong enough not
merely to constrain a spectrum but to \emph{produce} one. Given a gauge
algebra and the divisor class it sits on, the multiplicities of the charged
matter are the unique solution of three linear equations, in exact rational
arithmetic on the intersection form of a surface. There is no cohomology
computation, no metric, no approximation, and --- this is the part worth
insisting on --- no table.

\Cref{sec:kodaira,sec:anomaly,sec:spinors,sec:clusters} develop that machinery
and check it against itself. \Cref{sec:bases,sec:chi,sec:mw} turn to the base
surface, where several quantities admit two unrelated computations that must
agree. \Cref{sec:four} goes up one dimension to the elliptic fourfold and the
$G$-flux, where the exact statements stop early and the paper says so.
\Cref{sec:cicy} connects the constructions back to the complete intersections
the companion paper is about. \Cref{sec:orientifold,sec:sen} treat Type IIB
orientifolds and Sen's weak coupling limit, which is where the two halves of
the paper meet. \Cref{sec:boundary} restates the boundary.

Everything is implemented in \texttt{pyCICY.theories}, which is the same
subpackage the heterotic constructions of~\cite{SuppMath} live in and shares
its interface; every number quoted below is recomputed by
\texttt{paper/make\_strings\_facts.py} and inserted into this document
mechanically.

% =====================================================================

\section{Kodaira's table, and its internal consistency}
\label{sec:kodaira}

An F-theory background is a Weierstrass model
%
\begin{equation}
y^2 \;=\; x^3 + f\,x + g ,
\qquad
f \in H^0(B, -4K_B),
\quad
g \in H^0(B, -6K_B),
\label{eq:weierstrass}
\end{equation}
%
over a base $B$, with the fibre degenerating over the vanishing locus of the
discriminant $\Delta = 4f^3 + 27g^2 \in |-12K_B|$. The type of degeneration
--- hence the gauge algebra carried by that divisor --- is fixed by the orders
to which $f$, $g$ and $\Delta$ vanish along it, by Kodaira's
classification~\cite{Kodaira1963,BIKMSV1996}.

The table is quoted, as it must be. What is not quoted is its consistency with
the equation it classifies. Since $\Delta = 4f^3 + 27g^2$,
%
\begin{equation}
\operatorname{ord}\Delta \;\geq\; \min\bigl(3\operatorname{ord}f,\;
2\operatorname{ord}g\bigr),
\label{eq:deltabound}
\end{equation}
%
with equality unless the two terms cancel at leading order, which they cannot
when $3\operatorname{ord}f \neq 2\operatorname{ord}g$. This is an arithmetic
constraint independent of Kodaira, and a classifier that assigns a fibre type
to orders violating it is assigning physics to a geometry that does not exist.
The implementation enforces \cref{eq:deltabound}, and a scan over all
$(\operatorname{ord}f, \operatorname{ord}g, \operatorname{ord}\Delta)$ in a
box confirms that no entry of the table is rejected by it. That is a real
check: the $I_0^*$ row in particular requires
$\operatorname{ord}f = 2$ \emph{or} $\operatorname{ord}g = 3$, not merely both
at least those values, and without that refinement the classifier accepts
$(3,4,6)$, which \cref{eq:deltabound} forbids.

A model with $\operatorname{ord}f \geq 4$ and $\operatorname{ord}g \geq 6$ is
non-minimal: it is not the Weierstrass form of a Calabi--Yau with canonical
singularities, and the base must be blown up before it means anything. This is
reported as such rather than being given a gauge algebra, and the distinction
matters in \cref{sec:bases}, where it disqualifies three of the Hirzebruch
surfaces as bases.

% =====================================================================

\section{Anomaly cancellation as a source of spectra}
\label{sec:anomaly}

\subsection{The conditions}

For a simple algebra $\g{g}$ on an irreducible divisor $D$ in the base, write
$A_R$, $B_R$, $C_R$ for the coefficients in
%
\begin{equation}
\operatorname{tr}_R F^2 = A_R \operatorname{tr}F^2 ,
\qquad
\operatorname{tr}_R F^4 = B_R \operatorname{tr}F^4
+ C_R \bigl(\operatorname{tr}F^2\bigr)^2 ,
\label{eq:traces}
\end{equation}
%
the trace on the right being in the defining representation, and $\lambda$ for
the normalisation of~\cite{KumarMorrisonTaylor2010}. Anomaly cancellation
requires
%
\begin{align}
\lambda \Bigl( A_{\mathrm{adj}} - \sum_R x_R A_R \Bigr)
&\;=\; 6\,K \cdot D , \label{eq:condA}\\
\lambda^2 \Bigl( \sum_R x_R C_R - C_{\mathrm{adj}} \Bigr)
&\;=\; 3\,D \cdot D , \label{eq:condC}\\
B_{\mathrm{adj}} &\;=\; \sum_R x_R B_R , \label{eq:condB}
\end{align}
%
the last only when $\g{g}$ has an independent quartic Casimir. The adjoint
appears on both sides of the first two: once as the vector multiplet, and once
as matter with multiplicity equal to the genus of $D$, since a curve of genus
$g$ carries $g$ adjoint hypermultiplets. The intersection $K \cdot D$ is not
independent data; adjunction fixes it as $2g - 2 - D\cdot D$. Alongside these
sits the gravitational condition
%
\begin{equation}
H - V + 29\,T \;=\; 273 ,
\label{eq:grav}
\end{equation}
%
with $T = h^{1,1}(B) - 1$ tensor multiplets.

\subsection{Matter, solved}
\label{sec:matter}

\Crefrange{eq:condA}{eq:condB} are three linear equations in the
multiplicities. For $\g{su}(N)$ with fundamental and antisymmetric matter they
are three equations in two unknowns and are consistent, which is not
automatic, and their solution reproduces the standard results without any of
them being stored:

\begin{center}
\begin{tabular}{lccc}
\toprule
$D\cdot D$ & $-2$ & $-1$ & $0$ \\
\midrule
fundamentals of $\g{su}(5)$
& \SFSUFiveMinusTwoFund & \SFSUFiveMinusOneFund & \SFSUFiveZeroFund \\
antisymmetrics
& --- & \SFSUFiveMinusOneAnti & \SFSUFiveZeroAnti \\
\bottomrule
\end{tabular}
\end{center}

\noindent
that is, $2N$ fundamentals on a $-2$ curve, $N+8$ and one antisymmetric on a
$-1$ curve, sixteen and two on a $0$ curve.

The case worth dwelling on is $\g{e}_7$ on a curve of self-intersection $-7$,
where the unique solution is $x = \SFESevenHalf$ of a $\mathbf{56}$. Half a
hypermultiplet is not a rounding artefact: the $\mathbf{56}$ is pseudo-real,
the reality condition can be imposed on half of it, and the answer is exactly
one half. The rational arithmetic is carried through with
\texttt{fractions.Fraction} for precisely this reason. The
\SFESevenCharged{} charged states it contributes are what
\cref{eq:grav} then sees.

\subsection{The matter-free algebras, derived}
\label{sec:matterfree}

Set every multiplicity to zero on a rational curve and
\cref{eq:condA,eq:condC} each become a formula for the self-intersection:
%
\begin{equation}
D \cdot D \;=\; -2 - \frac{\lambda A_{\mathrm{adj}}}{6}
\qquad\text{and}\qquad
D \cdot D \;=\; -\frac{\lambda^2 C_{\mathrm{adj}}}{3} .
\label{eq:matterfree}
\end{equation}
%
These are different functions of the group theory. They agree for very few
algebras and at an integer for fewer still, and the survivors are exactly
\SFMatterFreeCount{} in number:

\begin{center}
\begin{tabular}{lcccccc}
\toprule
algebra & $\g{su}(3)$ & $\g{so}(8)$ & $\g{f}_4$ & $\g{e}_6$ & $\g{e}_7$
& $\g{e}_8$ \\
\midrule
$D\cdot D$ & $-\SFMFsuThree$ & $-\SFMFsoEight$ & $-\SFMFfFour$
& $-\SFMFeSix$ & $-\SFMFeSeven$ & $-\SFMFeEight$ \\
\bottomrule
\end{tabular}
\end{center}

\noindent
which is the tabulated list of matter-free non-Higgsable
clusters~\cite{MorrisonTaylor2012}, at the tabulated self-intersections.
Nothing in the derivation consults that list. $\g{g}_2$ and $\g{su}(2)$ fail
the integrality of \cref{eq:matterfree} and are therefore forced to carry
matter, which is also correct and is what \cref{sec:clusters} exploits.

The single-curve non-Higgsable table is quoted, since which algebra the
\emph{generic} Weierstrass model realises on a curve is not an anomaly
question. But its matter content is derived, and the agreement of the derived
matter-free subset with the quoted table is a check on the group theory in
\cref{eq:traces} rather than an input to it.

% =====================================================================

\section{The \texorpdfstring{$\g{so}(N)$}{so(N)} spinors, from their weights}
\label{sec:spinors}

The vector and adjoint coefficients of $\g{so}(N)$ are elementary. The spinor
coefficients are usually taken from a table; here they are computed. The
weights of a spinor are $(\pm\tfrac12, \dots, \pm\tfrac12)$ in the orthogonal
basis, so placing the field strength in the Cartan with eigenvalues
$x_1, \dots, x_r$ and summing over sign patterns gives the traces directly.
The odd terms drop out of the sign sum; the cross terms drop out because
summing a product of an incomplete set of signs over all patterns vanishes.
Writing $t_2$ and $t_4$ for the vector traces,
%
\begin{equation}
\operatorname{tr}_S F^2 = \frac{d_S}{8}\, t_2 ,
\qquad
\operatorname{tr}_S F^4 = -\frac{d_S}{16}\, t_4
+ \frac{3 d_S}{64}\, t_2^2 ,
\label{eq:spinor}
\end{equation}
%
with $d_S$ the dimension of the spinor. The even and odd rank cases have
different weight lattices and different numbers of sign patterns, and they
land on the same formula once written in terms of $d_S$; that they do is the
check that the two were done consistently. For $\g{so}(10)$,
\cref{eq:spinor} gives the $\mathbf{\SFSOTenSpinorDim}$ with
$A = \SFSOTenSpinorA$.

$\g{so}(8)$ is the exception, and the exception is informative. There the
quartic expansion reaches a term involving all four signs at once, which does
not cancel. That term is the Pfaffian --- the extra invariant $\g{so}(8)$
alone possesses --- and triality is the statement that its two spinors are
then indistinguishable from the vector. They take the vector's coefficients,
and consequently the spinor column duplicates the vector column and would
leave the anomaly system underdetermined; the implementation trims it from the
defaults and says why.

With spinors available, $\g{so}(N)$ is no longer confined to the $-4$ curve
where the vector alone would put it. On a $-3$ curve $\g{so}(10)$ carries
\SFSOTenMinusThreeVec{} vectors and \SFSOTenMinusThreeSpin{} spinor, and
\crefrange{eq:condA}{eq:condB} remain three equations in two unknowns and
remain consistent.

\subsection{Reality, and which multiplicities exist}

A half-hypermultiplet exists only for a pseudo-real representation, and the
$\g{so}(N)$ spinors follow the eightfold pattern: real for
$N \equiv 0, 1, 7 \bmod 8$, pseudo-real for $3, 4, 5$, complex for $2, 6$.
This is not bookkeeping. On a $-3$ curve $\g{so}(12)$ is given
$\SFSOTwelveHalfSpinor$ of a $\mathbf{32}$, which is legitimate because that
spinor is pseudo-real; $\g{so}(14)$ on the same curve is given a quarter of a
$\mathbf{64}$, which is not a spectrum at all, and $\g{so}(16)$ a fraction of
a real representation, which is likewise not. Both are refused. A
multiplicity that is neither an integer nor --- in the pseudo-real case --- a
half is a statement that the configuration does not exist, and it is more
useful reported as such than rounded.

% =====================================================================

\section{Clusters spanning several curves}
\label{sec:clusters}

A curve too shallow to force a gauge algebra on its own can belong to a chain
that forces one collectively, because the intersections tie the Weierstrass
model on one curve to the model on its neighbour. There are exactly
\SFClusterCount{} such clusters~\cite{MorrisonTaylor2012}, and their matter is
derived here from the same conditions, extended by the mixed one.

\subsection{A normalisation the clusters fix}

For two simple factors the mixed anomaly condition is
%
\begin{equation}
b_i \cdot b_j \;=\; \lambda_i \lambda_j
\sum_{R,S} x^{ij}_{RS}\, A_R\, A_S .
\label{eq:mixed}
\end{equation}
%
Every representation of index one --- the fundamental of $\g{su}$ and
$\g{sp}$, the vector of $\g{so}$, the $\mathbf{27}$, $\mathbf{56}$,
$\mathbf{26}$ and $\mathbf{7}$ --- makes the sum on the right the
multiplicity, so that for two $\g{su}$ factors the intersection number
\emph{is} the number of bifundamentals. The normalisation factors do not
cancel in general, and the $(-3,-2)$ cluster is what settles them. There
$\g{g}_2$ sits on the $-3$ curve and $\g{su}(2)$ on the $-2$, meeting once.
Their own conditions give the $\g{g}_2$ exactly $\SFGTwoSevens$ seven and the
$\g{su}(2)$ exactly $\SFGTwoSUTwoDoublets$ doublets. A full
$(\mathbf{7},\mathbf{2})$ would require two sevens, which are not there;
\cref{eq:mixed} with the $\lambda$'s gives multiplicity
$\SFGTwoSharedMult$, which requires exactly the one that is. Dropping the
$\lambda$'s is invisible when both factors are $\g{su}$ and wrong as soon as
one is not.

\subsection{The three-curve cluster}

The $(-2,-3,-2)$ chain is the sharper case, because the middle curve's matter
is claimed from both sides. Its own conditions give $\g{so}(7)$ on the $-3$
curve $\SFSOSevenSpinors$ spinors and no vectors at all, and each $\g{su}(2)$
neighbour claims one of them through a half
$(\mathbf{2},\mathbf{8})$; the accounting closes exactly, with
$\SFClusterThreeCharged$ charged states in the cluster.

Two things had to be right for that to happen. First, the spinor coefficients
of \cref{eq:spinor} for $\g{so}(7)$, which are not the tabulated ones for any
larger $\g{so}$ and which this cluster tests independently. Second, the
identity of the representation carrying the shared matter: $\g{so}(7)$ has
\emph{two} representations of index one, namely \SFSOSevenIndexOne, and on a
$-3$ curve it carries no vectors, so the bifundamental sits in the spinor.
Which one appears is decided by the matter the curve actually has and not by
the algebra, and assuming the vector --- the natural default --- looks for
states that are not there.

% =====================================================================

\section{The base surface, counted twice}
\label{sec:bases}

Six-dimensional F-theory needs remarkably little about the base: the
intersection pairing on $H^2$ and the canonical class. Two computations of the
same quantity out of that data, sharing nothing, are available and must agree.

The complex structure moduli of the generic Weierstrass model are a
Riemann--Roch count. The model is fixed by $f$ and $g$ of \cref{eq:weierstrass}
modulo the automorphisms of $B$ and the rescaling
$(f,g) \mapsto (t^4 f, t^6 g)$, so
%
\begin{equation}
n_{\text{moduli}} \;=\; h^0(-4K) + h^0(-6K) - \chi(T_B) - 1 .
\label{eq:moduli}
\end{equation}
%
On $\PP^2$ that is $\SFPTwoHfour + \SFPTwoHsix - \SFPTwoChiT - 1 =
\SFPTwoModuli$. The gravitational anomaly \cref{eq:grav} is a one-loop
condition in six dimensions and says the answer must be $272 - 29T$. The two
have nothing to do with each other, and they agree on every rational base
tested (\SFModuliAgree).

The delicate case is $F_2$. Its automorphism group is one dimension larger
than $F_0$'s, which on its own would undercount the moduli by one; the
compensation is that $F_2$ deforms to $F_0$, so $h^1(T_B) = 1$. Using the
Euler characteristic $\chi(T_B)$ in \cref{eq:moduli} rather than $h^0(T_B)$
alone is exactly what holds the count at $\SFFTwoModuli$ across the Hirzebruch
surfaces, with $\chi(T_B) = \SFFTwoChiT$ throughout. A formula in $h^0$ would
agree on $F_0$ and $F_1$ and fail on $F_2$, which is the kind of error that
survives casual testing.

\subsection{Whole models}

With the matter derived and \cref{eq:grav} closing the neutral count, the
Hodge numbers of the elliptic threefold follow from
$h^{1,1}(X) = h^{1,1}(B) + 1 + \operatorname{rank}(G)$ and
$h^{2,1}(X) = H_{\text{neutral}} - 1$:

\begin{center}
\begin{tabular}{lcccc}
\toprule
base & algebra & $h^{1,1}$ & $h^{2,1}$ & $\chi$ \\
\midrule
$\PP^2$   & \SFPTwoAlg    & \SFPTwoHone    & \SFPTwoHtwo    & \SFPTwoChi \\
$F_0$     & \SFFZeroAlg   & \SFFZeroHone   & \SFFZeroHtwo   & \SFFZeroChi \\
$F_7$     & $\g{e}_7$     & \SFFSevenHone  & \SFFSevenHtwo  & \SFFSevenChi \\
$F_{12}$  & $\g{e}_8$     & \SFFTwelveHone & \SFFTwelveHtwo & \SFFTwelveChi \\
\bottomrule
\end{tabular}
\end{center}

\noindent
$F_7$ is the only one of these with charged matter, and it is the
half-hypermultiplet of \cref{sec:matter}: \SFFSevenCharged{} charged states.
$F_{12}$ is the far end of the six-dimensional landscape and reproduces the
familiar $(\SFFTwelveHone, \SFFTwelveHtwo)$.

The bases $F_9$, $F_{10}$ and $F_{11}$ are absent on purpose. Their sections
carry points where the Weierstrass model is non-minimal in the sense of
\cref{sec:kodaira}, so they are not bases at all until those points are blown
up, after which the surface is no longer Hirzebruch. Reporting a spectrum for
them would be reporting physics for a geometry that does not exist. On the
heterotic side those are exactly the cases with three, two and one instanton
in the second $E_8$ --- too few to fit --- and the small instanton transition
that repairs it is the same event as the blowup~\cite{MorrisonVafaII}.

% =====================================================================

\section{Three routes to one Euler characteristic}
\label{sec:chi}

The Hodge numbers above give $\chi = 2(h^{1,1} - h^{2,1})$ through the
spectrum and hence through the anomaly. Independently, for a smooth
Weierstrass model over a surface, the Chern classes of the fibration
give~\cite{KLRY1998,GrassiMorrison2011}
%
\begin{equation}
\chi(X) \;=\; -60 \int_B c_1(B)^2 \;=\; -60\,K_B^2 .
\label{eq:chithree}
\end{equation}
%
On $\PP^2$ both give $\SFPTwoChern$. They agree on every base for which the
generic Weierstrass model is smooth, and they must \emph{disagree} where it is
not: $F_{12}$ gives $\SFFTwelveChi$ through the spectrum against
$\SFFTwelveChern$ through \cref{eq:chithree}, and the difference of
$\SFFTwelveGap$ is the resolution of the $\g{e}_8$ fibre. An implementation
reporting agreement there would not be computing what it claims. Together with
the moduli count of \cref{sec:bases} this makes three independent handles on
the same geometry, and the pattern of where they agree is as informative as
the agreement.

% =====================================================================

\section{Extra sections and the Mordell--Weil rank}
\label{sec:mw}

An elliptic fibration with rational sections beyond the zero section carries
abelian gauge factors~\cite{MorrisonPark2012}, one per generator of the
Mordell--Weil group, and each contributes a vector multiplet to $V$ and a
divisor to $h^{1,1}$. The abelian anomaly conditions are
\cref{eq:condA,eq:condC} with the adjoint removed and $\lambda A_R$ replaced
by $q^2$:
%
\begin{equation}
\sum_q x_q\, q^2 \;=\; 6\, c_1 \cdot b ,
\qquad
\sum_q x_q\, q^4 \;=\; 3\, b \cdot b ,
\label{eq:abelian}
\end{equation}
%
with $b$ the height pairing of the generating section, playing the role the
gauge divisor plays for a non-abelian factor.

Two conditions determine two charges and no more, and the implementation says
so rather than guessing. Setting every multiplicity to zero forces
$c_1 \cdot b = b \cdot b = 0$: a $U(1)$ with no charged matter has vanishing
height pairing, which is to say it is not there. This is the abelian
counterpart of \cref{sec:matterfree}, and it finds nothing, which is the
correct answer.

Over $\PP^2$ with charges $1$ and $2$, only the height pairings
$b \in \{\SFMWGood\}$ times the hyperplane admit an integral spectrum at all;
the rest want fractional numbers of states and are refused. At $b = 6H$ the
solution is \SFMWCharge{} states of charge one and none of charge two, giving
$(h^{1,1}, h^{2,1}) = (\SFMWHone, \SFMWHtwo)$ against the sectionless
$(\SFPTwoHone, \SFPTwoHtwo)$ --- the extra section visible in both entries, as
one more divisor and as charged states removed from the neutral count.

Abelian and non-abelian factors together are declined rather than
approximated. A state charged under both would be counted once by its gauge
divisor and once by the $U(1)$, and which states those are is not determined
by the data available here.

% =====================================================================

\section{Four dimensions, and the box the flux lives in}
\label{sec:four}

An elliptic Calabi--Yau fourfold over a threefold base gives four-dimensional
$\mathcal{N}=1$. Two computations of its Euler characteristic are again
available and again share nothing. The complex structure moduli are the
generalisation of \cref{eq:moduli}, and a Calabi--Yau fourfold
satisfies~\cite{KLRY1998}
%
\begin{equation}
\chi \;=\; 6\bigl(8 + h^{1,1} + h^{3,1} - h^{2,1}\bigr) ,
\label{eq:chifour}
\end{equation}
%
while separately, for the smooth Weierstrass model,
%
\begin{equation}
\chi(X) \;=\; 12 \int_B c_1 c_2 \;+\; 360 \int_B c_1^3 .
\label{eq:chifourchern}
\end{equation}

\begin{center}
\begin{tabular}{lcccccc}
\toprule
base & $\int c_1 c_2$ & $\int c_1^3$ & $h^{1,1}$ & $h^{3,1}$
& $\chi$ \eqref{eq:chifour} & $\chi$ \eqref{eq:chifourchern} \\
\midrule
$\PP^3$
& \SFPThreeCoCt & \SFPThreeCoCube & \SFPThreeHone & \SFPThreeHthree
& \SFPThreeChi & \SFPThreeChiChern \\
$\PP^1 \times \PP^2$
& \SFPOneTwoCoCt & \SFPOneTwoCoCube & \SFPOneTwoHone & \SFPOneTwoHthree
& \SFPOneTwoChi & \SFPOneTwoChiChern \\
\bottomrule
\end{tabular}
\end{center}

\noindent
The $\PP^3$ row is the standard example and both of its entries are the
familiar numbers. The D3-brane tadpole is then
$\chi/24 = \SFPThreeTadpole$~\cite{SethiVafaWitten1996}.

That is where the exact statements stop, and the paper stops with them. The
$G$-flux is bounded but not determined. Witten quantisation requires
$G + c_2(X)/2$ to be integral~\cite{Witten1996flux}; the tadpole
$N_{D3} + \tfrac12 \int G \wedge G = \chi/24$ with $N_{D3} \geq 0$ bounds
$\int G \wedge G \leq \SFPThreeFluxBound$ on $\PP^3$; and supersymmetry
demands that $G$ be primitive and of type $(2,2)$. None of this is a spectrum.
The chiral index is an index twisted by $G$, and $G$ is a choice this package
does not carry, so the spectrum is declined rather than returned.

% =====================================================================

\section{The complete intersections, and what they are not}
\label{sec:cicy}

The constructions above are stated in terms of a base and its intersection
form, and none of them needs the total space as a complete intersection.
That is not an accident of presentation: the generic Weierstrass model is a
hypersurface in a $\PP^{2,3,1}$ bundle and is not a CICY. But complete
intersections do fibre in elliptic curves, and the criterion is
combinatorial~\cite{AndersonGaoGrayLee2017}. Reordering a configuration matrix
into
%
\begin{equation}
\begin{pmatrix} F & 0 \\ C & B \end{pmatrix}
\label{eq:fibblock}
\end{equation}
%
with the top rows vanishing outside the left columns makes the top block a
complete intersection in its own ambient factors; when it has dimension one it
is a Calabi--Yau one-fold, an elliptic curve, and the manifold fibres over the
bottom block. The Calabi--Yau condition on the fibre is automatic and worth
seeing why: each row of the full matrix sums to its ambient dimension plus
one, and a fibre row is zero on every base column, so its entries in the fibre
columns already sum correctly.

Implemented from that definition and swept over the published list, this finds
\SFCicyElliptic{} of the \SFCicyTotal{} CICY threefolds with an obvious
genus-one fibration, which is the published count, and \SFCicyKThree{} with an
obvious K3 fibration.

The gap between that and an F-theory background is the point of stating it. An
obvious fibration is a \emph{genus-one} fibration: the fibres are elliptic
curves but nothing marks a point on them, and without a section there is no
Weierstrass form and no place to read the axio-dilaton. The $(3,3)$
hypersurface in $\PP^2 \times \PP^2$ fibres over $\PP^2$ in plane cubics
\SFBicubicFibrations{} ways, one per projection, and has
$(h^{1,1}, h^{2,1}) = (\SFBicubicHone, \SFBicubicHtwo)$; the Weierstrass model
over the same base has $(\SFPTwoHone, \SFPTwoHtwo)$. Same base, same fibre
type, different manifolds. The tetraquadric fibres \SFTetraFibrations{} ways,
one per choice of two of its four $\PP^1$ factors.

% =====================================================================

\section{Type IIB orientifolds}
\label{sec:orientifold}

An orientifold quotients Type IIB by $\Omega_p (-1)^{F_L} \sigma$ with
$\sigma$ a holomorphic involution of $X$, and the four-dimensional closed
string spectrum follows from how $\sigma$ splits the cohomology. Here the
construction does return to complete intersections, because the involutions
that matter act on the ambient coordinates.

\subsection{Two kinds, distinguished by geometry}

A holomorphic involution acts on the holomorphic three-form by a sign, and the
sign decides everything: $\sigma^*\Omega = -\Omega$ gives O3- and O7-planes,
fixed points and divisors; $\sigma^*\Omega = +\Omega$ gives O5- and O9-planes,
fixed curves and all of $X$. The implementation computes the sign by counting
flipped coordinates and polynomial signs, and computes the fixed dimensions by
intersecting the defining polynomials with coordinate subspaces. That even
complex codimension accompanies the minus sign is a theorem, so the agreement
of the two is a test and not a definition.

The polynomial signs are not decoration. Flipping a set of coordinates and
flipping its complement are the \emph{same} map on projective space, and the
two descriptions give opposite values of $\sum_i \mathrm{minus}_i$; what puts
them back together is that they also give opposite $\varepsilon_a$ in
$p_a \circ \sigma = \varepsilon_a p_a$. Omitting them would make the O-plane
type depend on how the same involution was written down.

\subsection{The equivariant Hodge numbers, twice}
\label{sec:hodgesplit}

The spectrum is a function of the splits
$h^{1,1} = h^{1,1}_+ + h^{1,1}_-$ and $h^{2,1} = h^{2,1}_+ + h^{2,1}_-$, and
these are computed by two routes with no code in common.

The first is the topological Lefschetz fixed point theorem. For a holomorphic
involution of a Calabi--Yau threefold,
%
\begin{equation}
\chi(\mathrm{Fix}\,\sigma) \;=\;
2 + 2\bigl(h^{1,1}_+ - h^{1,1}_-\bigr) - 2 s
- 2 \bigl(h^{2,1}_+ - h^{2,1}_-\bigr) ,
\label{eq:lefschetz}
\end{equation}
%
with $s$ the sign on $\Omega$, using that the traces on $H^2$ and $H^4$ agree
by Poincar\'e duality. On a favourable CICY a coordinate sign flip fixes every
ambient hyperplane class, so $h^{1,1}$ is entirely invariant and
$\chi(\mathrm{Fix}\,\sigma)$ determines the rest. The fixed locus is a
complete intersection in a product of projective spaces and is almost never
Calabi--Yau, so it is computed from Chern classes rather than by the CICY
machinery.

The second route counts monomials: the complex structure deformations of a
hypersurface are its monomials modulo the ambient reparametrisations, and the
involution grades both.

On the quintic with one coordinate flipped, the fixed locus is a quintic
surface --- the O7 --- together with the point $[0{:}0{:}0{:}0{:}1]$, which
lies on $X$ because no invariant quintic has an $x_4^5$ term. Its Euler
characteristic is $\SFQuinticFixedChi$, and \cref{eq:lefschetz} gives
$(h^{2,1}_+, h^{2,1}_-) = (\SFQuinticHtwoPlus, \SFQuinticHtwoMinus)$. The
monomial count gives $(\SFQuinticDefPlus, \SFQuinticDefMinus)$ --- the same
numbers in the opposite order --- and reconciling them is the subtlety this
section exists for.

\begin{framed}
\noindent
$H^{2,1} \cong H^1(X, T_X) \otimes H^{3,0}$, and the monomials grade
$H^1(T_X)$, not $H^{2,1}$. When $\sigma^*\Omega = -\Omega$ the two gradings
are opposite, so the invariant deformations are the \emph{anti}-invariant part
of $H^{2,1}$. Getting this backwards leaves both routes internally consistent
and both wrong; the fixed point theorem is what decides which way round it
goes. With the twist applied the two agree (\SFQuinticRoutesAgree).
\end{framed}

\subsection{Involutions that change the manifold}

Some sign patterns cannot be imposed on a generic $X$ at all, and an
implementation that does not notice will return numbers belonging to a
different manifold. If every defining polynomial vanishes identically on a
coordinate subspace, that subspace lies inside $X$. When it is a divisor,
$X$ contains a linear subspace of the ambient whose class is not in the
lattice the ambient generates: $h^{1,1}$ exceeds what the configuration matrix
says, and the CICY Hodge numbers no longer describe $X$. When it has the
dimension of $X$, the polynomial factorises and $X$ is not irreducible.

On the quintic both occur. Flipping three coordinates forces every invariant
monomial to vanish on the plane they span, so $X$ contains a plane --- and a
quintic containing a plane is a Noether--Lefschetz jump with $h^{1,1} = 2$,
not $1$. Flipping four forces every invariant monomial to be divisible by the
fifth coordinate. Both routes of \cref{sec:hodgesplit} refuse in both cases,
and refuse for the stated reason rather than by failing an assertion. A curve
inside $X$ is not a problem, since a generic quintic contains lines; only a
divisor is.

\subsection{Both projections}

For $\sigma^*\Omega = -\Omega$ the massless closed string sector is
$h^{1,1}_+$ chiral multiplets pairing the K\"ahler moduli with $C_4$,
$h^{1,1}_-$ from $B_2$ and $C_2$, $h^{2,1}_-$ complex structure moduli,
$h^{2,1}_+$ vectors from $C_4$, and the axio-dilaton~\cite{GrimmLouis2004}.
For $\sigma^*\Omega = +\Omega$ the projection is $\Omega_p \sigma$ without the
left-moving fermion number, so $C_0$, $C_2$ and $C_4$ all change parity ---
the statement that Type I retains $C_2$. Running the reduction again:

\begin{center}
\begin{tabular}{lll}
\toprule
& O3/O7 & O5/O9 \\
\midrule
K\"ahler moduli        & $h^{1,1}_+$ & $h^{1,1}_+$ \\
two-form moduli        & $h^{1,1}_-$ & $h^{1,1}_-$ \\
complex structure      & $h^{2,1}_-$ & $h^{2,1}_+$ \\
vector multiplets      & $h^{2,1}_+$ & $h^{2,1}_-$ \\
remaining chiral       & axio-dilaton & dilaton with the dual of $C_{\mu\nu}$ \\
\bottomrule
\end{tabular}
\end{center}

\noindent
The K\"ahler moduli stay in $h^{1,1}_+$ either way, because the K\"ahler form
is invariant under an isometry regardless of the sign. What swaps is the
three-form sector, and it swaps for the reason in \cref{sec:hodgesplit}: the
deformations live in $H^1(T_X)$ and $\sigma$ acts on $H^{3,0}$ by the
$\Omega$ sign, so which half of $H^{2,1}$ is matter and which is gauge follows
the sign. There is no axio-dilaton in the second column: $C_0$ is odd under
worldsheet parity and there is no $\sigma$-odd zero-form for it to survive on.
One chiral multiplet either way, but not the same one.

On the quintic the two involutions above give
$\SFQuinticChiral$ chiral and $\SFQuinticVectors$ vector multiplets in the
O3/O7 case, and $\SFQuinticFiveChiral$ chiral and $\SFQuinticFiveVectors$
vector in the O5/O9 case, from the split
$(\SFQuinticFiveHtwoPlus, \SFQuinticFiveHtwoMinus)$.

The open string sector is absent throughout. The gauge group is a choice of
D7-brane divisor classes subject to
$\sum_a N_a ([D_a] + [D_a']) = 8[O7]$, not a consequence of the involution,
and the condition is stated rather than solved.

% =====================================================================

\section{Sen's limit, where the two halves meet}
\label{sec:sen}

Tuning the Weierstrass coefficients to
$f = -3h^2 + \epsilon\,\eta$ and $g = -2h^3 + \epsilon\,h\eta -
\epsilon^2 \psi$ and taking $\epsilon \to 0$ drives the axio-dilaton to weak
coupling almost everywhere~\cite{Sen1997}. The discriminant degenerates to
$\epsilon^2 h^2 (\eta^2 - h\psi)$, a double zero on the O7-plane $\{h = 0\}$
and a single brane on the rest, and the physical space is the double cover of
$B$ branched over the O7. Since $f$ is a section of $-4K_B$ and $h^2$ must be
too, $[O7] = -2K_B$, and everything follows from that one fact.

\begin{enumerate}[leftmargin=2em]
\item \emph{The cover is always a K3.} A double cover branched over $C$ has
$\chi = 2\chi(B) - \chi(C)$, and for a rational base $\chi(B) = 3 + T$ while
$\chi(C) = -2(9-T)$, so
%
\begin{equation}
\chi \;=\; 2(3+T) + 2(9-T) \;=\; 24 ,
\label{eq:k3}
\end{equation}
%
the tensor multiplet count cancelling because $K^2 = 9 - T$ on a rational
surface. Over the \SFSenBases{} bases tested the cover has
$\chi \in \{\SFSenCoverChi\}$. The genus of the branch curve does vary ---
\SFSenGenusPTwo{} on $\PP^2$, \SFSenGenusDPEight{} on $dP_8$ --- but the cover
does not.

\item \emph{The D7 tadpole closes for every base at once.} The brane at
$\eta^2 = h\psi$ lies in $|-8K|$, and with its orientifold image that is
$-16K = 8[O7]$. Both sides are multiples of the same canonical class
(\SFSenTadpoleAll).

\item \emph{The brane rules reproduce Kodaira.} This is the sharpest check
in the package.
\end{enumerate}

Perturbatively a stack of $n$ D7-branes gives $\g{u}(n)$, or $\g{so}(2n)$ when
it sits on the O7-plane and its images coincide with it --- Chan--Paton
bookkeeping and nothing more. Non-perturbatively the same configuration is a
degeneration of the elliptic fibration with a fibre type read off the
vanishing orders of $f$, $g$ and $\Delta$:

\begin{center}
\begin{tabular}{lll}
\toprule
configuration & Chan--Paton & Kodaira \\
\midrule
$n$ branes, away from the O7 & $\g{u}(n)$    & $I_n$ carries $\g{su}(n)$ \\
four branes, on the O7       & $\g{so}(8)$   & $I_0^*$ carries $\g{so}(8)$ \\
$n$ branes, on the O7        & $\g{so}(2n)$  & $I_{n-4}^*$ carries $\g{so}(2n)$ \\
\bottomrule
\end{tabular}
\end{center}

\noindent
The shift by four is the content: it takes four D7-branes to cancel the charge
of an O7-plane locally, which is why $n$ branes on it give $I_{n-4}^*$ rather
than $I_n^*$, and why four give the smooth $I_0^*$. One side is a count of
Chan--Paton factors and the other is the vanishing order of a discriminant;
neither computation knows about the other, and they agree
(\SFBraneKodairaAgree).

% =====================================================================

\section{The boundary, restated}
\label{sec:boundary}

The companion paper distinguishes three ways a quantity can fail to be
available~\cite{SuppMath}. These constructions need a fourth, and it belongs
at the head of the list rather than the tail, because it is not a failure.

\begin{enumerate}[leftmargin=2em]
\item \emph{Does not exist.} The Yukawa couplings of a six-dimensional
F-theory compactification. $\mathcal{N}=(1,0)$ supersymmetry forbids a
superpotential: hypermultiplets are quaternionic and there is no
gauge-invariant holomorphic cubic coupling to write down. Returning zero would
be numerically right and conceptually wrong, and blaming the metric would
blame it for an absence it has nothing to do with. The implementation raises a
distinct exception, \texttt{NoSuchTheory}, and the tests require that it is
\emph{not} the metric exception.

\item \emph{Not implemented, but exact in principle.} The Yukawa couplings of
a four-dimensional orientifold, which come from triple overlaps of open string
wavefunctions on intersecting D7-brane divisors together with worldsheet
instantons. These exist; they are simply not computed. The distinction from
the previous item is the whole point of keeping the two exceptions apart.

\item \emph{Not determined by the data carried.} The chiral spectrum of a
four-dimensional F-theory model, which is an index twisted by the $G$-flux;
the decomposition of the D3 tadpole; the gauge group of an orientifold, which
is a choice of D7-brane classes subject to a stated condition. In each case
the constraints are exact and the choice is absent, and the constraints are
reported.

\item \emph{Obstructed by the metric.} The hypermultiplet moduli space metric
and the normalisation of the kinetic terms. Nothing in
\crefrange{sec:kodaira}{sec:sen} depends on these, which is why so much of the
above is exact.
\end{enumerate}

\noindent
The reason the third category is large here and small in~\cite{SuppMath} is
structural. A heterotic model on a CICY is specified by data the configuration
matrix and a bundle already contain; an F-theory or orientifold model is
specified by a base together with choices --- fluxes, brane stacks, sections
--- that live outside the geometry. What the anomaly conditions do is make
those choices finite and checkable rather than open-ended: a bound of
$\SFPThreeFluxBound$ on $\int G \wedge G$ is not a spectrum, but it is a box,
and a box is what an exhaustive method needs.

% =====================================================================

\section{Conclusion}
\label{sec:conclusion}

Where the companion paper found that a configuration matrix determines more
than one expects, this one finds that consistency does. The six-dimensional
anomaly conditions are usually presented as a filter applied to a spectrum
obtained some other way; used as equations rather than inequalities they
produce the spectrum outright, and they are restrictive enough that the
answers come out unique. The matter content of every gauge algebra on every
divisor, the six algebras that need no charged matter, the half-hypermultiplet
of $\g{e}_7$, the trace coefficients of the $\g{so}(N)$ spinors, the way the
multi-curve clusters share their matter: all of it is linear algebra over
$\mathbb{Q}$, and all of it is checkable against a table nobody had to type
in.

The recurring methodological point is that a quantity computed once is a
result and a quantity computed twice is a test. The moduli count of
\cref{sec:bases}, the Euler characteristics of \cref{sec:chi,sec:four}, the
equivariant Hodge numbers of \cref{sec:hodgesplit} and the brane rules of
\cref{sec:sen} are each obtained by two routes that share no code and,
usually, no branch of mathematics. Every one of the errors found during this
work was found by one of those pairs disagreeing --- a normalisation dropped
from \cref{eq:mixed}, a twist by $\Omega$ applied the wrong way, a fibre type
assigned to vanishing orders \cref{eq:deltabound} forbids --- and none of them
would have been visible to a single computation, however carefully written.

% =====================================================================

\section{Source code and limitations}
\label{sec:source}

Everything above is implemented in \texttt{pyCICY.theories.ftheory} and
\texttt{pyCICY.theories.orientifold}, with
\texttt{tests/test\_ftheory.py} and \texttt{tests/test\_orientifold.py}
executing the cross-checks quoted here. The numbers in this document are
emitted by \texttt{paper/make\_strings\_facts.py}. The repository is at
\url{https://github.com/brentharts/CICY}.

Known limitations, stated so that the scope is not mistaken:
%
\begin{itemize}[leftmargin=2em,itemsep=1pt]
\item Orientifold involutions are restricted to coordinate sign flips.
Involutions permuting ambient factors act non-trivially on $H^{1,1}$ and
permute the defining polynomials, and the restriction of a permuted polynomial
to the fixed locus depends on a sign the configuration matrix does not fix.
\item Multi-curve non-Higgsable clusters are handled locally, from curves and
their intersections; the automatic detection of forced gauge algebras on a
given base covers only the single-curve case.
\item Spinor coefficients are tabulated for $\g{so}(N)$ only; the $\g{sp}$ and
exceptional matter beyond the representations listed is not covered.
\item The Mordell--Weil sector handles the rank and the abelian conditions
with two distinct charges, and declines models mixing abelian and non-abelian
charged matter.
\item $G$-flux is bounded, not enumerated.
\end{itemize}

% =====================================================================

\begin{thebibliography}{99}
\setlength{\itemsep}{1pt plus 0.5pt}

\bibitem{SuppMath}
B.S.~Hartshorn,
\emph{What a configuration matrix determines, and what it does not: exact
methods for heterotic compactifications on CICYs},
companion document (Pre-print 2026),
\url{https://doi.org/10.5281/zenodo.21844007}.

\bibitem{SuppFigures}
B.S.~Hartshorn,
\emph{Supplementary material: figures for the pyCICY-X package},
companion document (Pre-print 2026),
\url{https://doi.org/10.5281/zenodo.21843383}.

\bibitem{Vafa1996}
C.~Vafa,
\emph{Evidence for F-theory},
Nucl.\ Phys.\ B \textbf{469} (1996) 403,
\href{https://arxiv.org/abs/hep-th/9602022}{arXiv:hep-th/9602022}.

\bibitem{MorrisonVafaI}
D.~R.~Morrison and C.~Vafa,
\emph{Compactifications of F-theory on Calabi--Yau threefolds I},
Nucl.\ Phys.\ B \textbf{473} (1996) 74,
\href{https://arxiv.org/abs/hep-th/9602114}{arXiv:hep-th/9602114}.

\bibitem{MorrisonVafaII}
D.~R.~Morrison and C.~Vafa,
\emph{Compactifications of F-theory on Calabi--Yau threefolds II},
Nucl.\ Phys.\ B \textbf{476} (1996) 437,
\href{https://arxiv.org/abs/hep-th/9603161}{arXiv:hep-th/9603161}.

\bibitem{Kodaira1963}
K.~Kodaira,
\emph{On compact analytic surfaces II},
Ann.\ Math.\ \textbf{77} (1963) 563.

\bibitem{BIKMSV1996}
M.~Bershadsky, K.~Intriligator, S.~Kachru, D.~R.~Morrison, V.~Sadov and
C.~Vafa,
\emph{Geometric singularities and enhanced gauge symmetries},
Nucl.\ Phys.\ B \textbf{481} (1996) 215,
\href{https://arxiv.org/abs/hep-th/9605200}{arXiv:hep-th/9605200}.

\bibitem{GSW1985}
M.~B.~Green, J.~H.~Schwarz and P.~C.~West,
\emph{Anomaly free chiral theories in six dimensions},
Nucl.\ Phys.\ B \textbf{254} (1985) 327.

\bibitem{Sagnotti1992}
A.~Sagnotti,
\emph{A note on the Green--Schwarz mechanism in open string theories},
Phys.\ Lett.\ B \textbf{294} (1992) 196,
\href{https://arxiv.org/abs/hep-th/9210127}{arXiv:hep-th/9210127}.

\bibitem{Erler1994}
J.~Erler,
\emph{Anomaly cancellation in six dimensions},
J.\ Math.\ Phys.\ \textbf{35} (1994) 1819,
\href{https://arxiv.org/abs/hep-th/9304104}{arXiv:hep-th/9304104}.

\bibitem{KumarMorrisonTaylor2010}
V.~Kumar, D.~R.~Morrison and W.~Taylor,
\emph{Global aspects of the space of 6D $\mathcal{N}=1$ supergravities},
JHEP \textbf{11} (2010) 118,
\href{https://arxiv.org/abs/1008.1062}{arXiv:1008.1062}.

\bibitem{MorrisonTaylor2012}
D.~R.~Morrison and W.~Taylor,
\emph{Classifying bases for 6D F-theory models},
Cent.\ Eur.\ J.\ Phys.\ \textbf{10} (2012) 1072,
\href{https://arxiv.org/abs/1201.1943}{arXiv:1201.1943}.

\bibitem{Taylor2011}
W.~Taylor,
\emph{TASI lectures on supergravity and string vacua in various dimensions},
\href{https://arxiv.org/abs/1104.2051}{arXiv:1104.2051}.

\bibitem{Weigand2018}
T.~Weigand,
\emph{F-theory},
PoS \textbf{TASI2017} (2018) 016,
\href{https://arxiv.org/abs/1806.01854}{arXiv:1806.01854}.

\bibitem{GrassiMorrison2011}
A.~Grassi and D.~R.~Morrison,
\emph{Anomalies and the Euler characteristic of elliptic Calabi--Yau
threefolds},
Commun.\ Num.\ Theor.\ Phys.\ \textbf{6} (2012) 51,
\href{https://arxiv.org/abs/1109.0042}{arXiv:1109.0042}.

\bibitem{KLRY1998}
A.~Klemm, B.~Lian, S.-S.~Roan and S.-T.~Yau,
\emph{Calabi--Yau fourfolds for M- and F-theory compactifications},
Nucl.\ Phys.\ B \textbf{518} (1998) 515,
\href{https://arxiv.org/abs/hep-th/9701023}{arXiv:hep-th/9701023}.

\bibitem{SethiVafaWitten1996}
S.~Sethi, C.~Vafa and E.~Witten,
\emph{Constraints on low-dimensional string compactifications},
Nucl.\ Phys.\ B \textbf{480} (1996) 213,
\href{https://arxiv.org/abs/hep-th/9606122}{arXiv:hep-th/9606122}.

\bibitem{Witten1996flux}
E.~Witten,
\emph{On flux quantization in M-theory and the effective action},
J.\ Geom.\ Phys.\ \textbf{22} (1997) 1,
\href{https://arxiv.org/abs/hep-th/9609122}{arXiv:hep-th/9609122}.

\bibitem{MorrisonPark2012}
D.~R.~Morrison and D.~S.~Park,
\emph{F-theory and the Mordell--Weil group of elliptically-fibered
Calabi--Yau threefolds},
JHEP \textbf{10} (2012) 128,
\href{https://arxiv.org/abs/1208.2695}{arXiv:1208.2695}.

\bibitem{Sen1997}
A.~Sen,
\emph{Orientifold limit of F-theory vacua},
Phys.\ Rev.\ D \textbf{55} (1997) R7345,
\href{https://arxiv.org/abs/hep-th/9702165}{arXiv:hep-th/9702165}.

\bibitem{GrimmLouis2004}
T.~W.~Grimm and J.~Louis,
\emph{The effective action of $\mathcal{N}=1$ Calabi--Yau orientifolds},
Nucl.\ Phys.\ B \textbf{699} (2004) 387,
\href{https://arxiv.org/abs/hep-th/0403067}{arXiv:hep-th/0403067}.

\bibitem{BKLS2007}
R.~Blumenhagen, B.~K\"ors, D.~L\"ust and S.~Stieberger,
\emph{Four-dimensional string compactifications with D-branes, orientifolds
and fluxes},
Phys.\ Rept.\ \textbf{445} (2007) 1,
\href{https://arxiv.org/abs/hep-th/0610327}{arXiv:hep-th/0610327}.

\bibitem{CollinucciDenefEsole2009}
A.~Collinucci, F.~Denef and M.~Esole,
\emph{D-brane deconstructions in IIB orientifolds},
JHEP \textbf{02} (2009) 005,
\href{https://arxiv.org/abs/0805.1573}{arXiv:0805.1573}.

\bibitem{AndersonGaoGrayLee2017}
L.~B.~Anderson, X.~Gao, J.~Gray and S.-J.~Lee,
\emph{Fibrations in CICY threefolds},
JHEP \textbf{10} (2017) 077,
\href{https://arxiv.org/abs/1708.07907}{arXiv:1708.07907}.

\bibitem{Candelas1988}
P.~Candelas, A.~M.~Dale, C.~A.~L\"utken and R.~Schimmrigk,
\emph{Complete intersection Calabi--Yau manifolds},
Nucl.\ Phys.\ B \textbf{298} (1988) 493.

\bibitem{Green1987}
P.~Green and T.~H\"ubsch,
\emph{Calabi--Yau manifolds as complete intersections in products of complex
projective spaces},
Commun.\ Math.\ Phys.\ \textbf{109} (1987) 99.

\bibitem{GrayHauptLukas2013}
J.~Gray, A.~S.~Haupt and A.~Lukas,
\emph{All complete intersection Calabi--Yau four-folds},
JHEP \textbf{07} (2013) 070,
\href{https://arxiv.org/abs/1303.1832}{arXiv:1303.1832}.

\bibitem{Larfors2022}
M.~Larfors, A.~Lukas, F.~Ruehle and R.~Schneider,
\emph{Learning size and shape of Calabi--Yau spaces},
\href{https://arxiv.org/abs/2111.01436}{arXiv:2111.01436};
code at \url{https://github.com/pythoncymetric/cymetric}.

\bibitem{Butbaia2024}
G.~Butbaia, D.~Mayorga Pe\~na, J.~Tan, P.~Berglund, T.~H\"ubsch,
V.~Jejjala and C.~Mishra,
\emph{Physical Yukawa couplings in heterotic string compactifications},
\href{https://arxiv.org/abs/2401.15078}{arXiv:2401.15078};
code at \url{https://github.com/Justin-Tan/cymyc}.

\end{thebibliography}

\end{document}
